\section{Introduction and precise statements}\label{sec:introduction}

Planar critical models are expected to have conformally invariant scaling
limits.  For random-cluster models this prediction identifies the law of an
interface, the joint law of all surrounding loops, and crossing
probabilities in arbitrary shapes.  These are distinct conclusions: a
limit for one interface does not by itself identify its coupling with
other interfaces, and convergence of loop traces on compact sets does not
by itself retain multiplicities, nesting, or complete traversals.

The random-cluster representation of Fortuin and Kasteleyn
\citep{FortuinKasteleyn1972} places percolation, Ising, and Potts models
in a common family with a real cluster weight.  Its planar loop and
transfer-matrix descriptions belong to the work of Temperley and Lieb
\citep{TemperleyLieb1971} and Baxter, Kelland, and Wu
\citep{BaxterKellandWu1976}.  Exact-solution and Coulomb-gas methods supplied
much of the prediction of critical behavior
\citep{Baxter1973,Nienhuis1984}.  We use that finite algebraic lineage,
and prove the analytic estimates needed to pass from the square-lattice
law to its conformal curve limits.

Schramm's characterization of conformally invariant interfaces led to
Schramm--Loewner evolution \citep{Schramm2000}.  Cardy's crossing prediction
\citep{Cardy1992} was proved by Smirnov for critical site percolation on
the triangular lattice \citep{Smirnov2001}; Camia and Newman subsequently
constructed its full continuum loop limit \citep{CamiaNewman2006Full}.
For FK-Ising, Smirnov's fermionic observable \citep{Smirnov2010} and the
isoradial analysis of Chelkak and Smirnov \citep{ChelkakSmirnov2012} led
to convergence of the Dobrushin interface to $\SLE_{16/3}$
\citep{ChelkakDCHonglerKemppainenSmirnov2014}.  Kemppainen and Smirnov
proved convergence of the boundary-touching FK-Ising loops
\citep{KemppainenSmirnov2019}, and then joint convergence of the full loop
ensemble and its exploration tree in bounded wired domains
\citep{KemppainenSmirnov2016FullScalingLimit}.

For general cluster weights, Duminil-Copin, Kozlowski, Krachun,
Manolescu, and Oulamara proved asymptotic rotational invariance of
critical square-lattice FK for $1\le q\le4$
\citep[Theorem~1.2]{DCKKMO2026Rotation}.
This establishes rotational invariance without identifying the limiting
interfaces or loop collections as SLE or CLE.  It is background for the
present result, whose identification proceeds through the four-change
observable.

There are also prior public claims for square-lattice bond percolation.
The preliminary preprint of Tsai, Yam, and Zhou
\citep{TsaiYamZhou2013SquareBond} claims $\SLE_6$ convergence of its
exploration path.  Zhou's later preprint \citep{Zhou2024SquareBond}
claims full-curve $\SLE_6$ convergence and Cardy convergence for its
stated crossing event under Condition C,
which imposes mesoscopic straight boundary segments and separation
conditions.  These are overlapping $q=1$ claims.  The statements here
include every uniformly approximated marked Jordan domain in the
specified lattice convention, every fixed $1\le q<4$, and the complete
nested whole-plane ensemble with multiplicities and traversal order.

For general cluster weights, the parafermionic program of Smirnov
\citep{Smirnov2010}, the work of Riva and Cardy \citep{RivaCardy2006},
and the criticality applications of Duminil-Copin
\citep{DuminilCopin2012Parafermion} establish the importance of local
winding identities.  Riva and Cardy explicitly distinguish such
identities from continuum holomorphicity, which also requires
regularity.  Chelkak and Smirnov's four-boundary-change construction
\citep[Section~6]{ChelkakSmirnov2012} supplies a further antecedent at
$q=2$.  The main analytic task here is to close this regularity gap for
the four-change observable throughout $1\le q<4$, and then to retain
the resulting tests after exploration.

The probabilistic inputs have similarly distinct roles.  Positive
association comes from the FKG inequality
\citep{FortuinKasteleynGinibre1971}.  Beffara and Duminil-Copin identify
the self-dual critical point \citep{BeffaraDuminilCopin2012};
Duminil-Copin, Sidoravicius, and Tassion prove continuity and buffered
crossing estimates \citep{DuminilCopinSidoraviciusTassion2017}.
We use the isoradial crossing estimates of Duminil-Copin, Li, and
Manolescu \citep{DCLM2019Isoradial} for transfer diagrams, and the
rough-quad and six-arm estimates of Duminil-Copin, Manolescu, and
Tassion \citep{DMT2021Fractal} for moving boundaries.  Their exact
parameter and boundary-condition ranges are stated below.

The continuum loop laws are built from branching SLE
\citep{Sheffield2009Exploration,MillerSheffield2016IG1,MillerSheffield2016IG3};
the simple-loop characterization \citep{SheffieldWerner2012} belongs
to the complementary simple regime.  For the nonsimple regime used
here, the large-domain construction and sphere invariance are supplied
by \citet{MillerWatsonWilson2015} and \citet{GwynneMillerQian2021}.
Our task is to identify the lattice tree and its complete loops with
these laws.  In doing so we distinguish convergence of drivers, ordered
curves, loop sets, and loop collections retaining multiplicity.
The results concern the nearest-neighbor square lattice in its ordinary
Euclidean embedding and its original random-cluster probability law.

\subsection{The model and its boundary conventions}
\label{intro:model}

Fix a real number $1\le q<4$ and put
\begin{equation}\label{intro:parameters}
 p_q=\frac{\sqrt q}{1+\sqrt q},
 \qquad
 \kappa(q)=\frac{4\pi}{\arccos(-\sqrt q/2)}.
\end{equation}
The inverse cosine has values in $[0,\pi]$.  Thus $4<\kappa(q)\le6$,
with $\kappa(1)=6$ and $\kappa(q)\downarrow4$ as $q\uparrow4$.  On a finite graph $G=(V,E)$,
let $\xi$ be a partition of designated boundary vertices.  For
$A\subseteq E$, let $k_\xi(A)$ be the number of components of $(V,A)$
after identification within each block of $\xi$, including isolated
vertices.  The random-cluster law is
\begin{equation}\label{intro:finite-law}
 \phi_{G,q}^{\xi}(A)
 =\frac{1}{Z_{G,q}^{\xi}}
 p_q^{|A|}(1-p_q)^{|E|-|A|}q^{k_\xi(A)}.
\end{equation}
An \emph{actual open path} always consists of open edges of $G$.
Identification in $\xi$ is used to compute the probability weight and
does not supply an additional edge in such a path.

Let $D$ be a bounded Jordan domain.  A lattice approximation at mesh
$\delta_n\downarrow0$ is the bounded interior $D_n$ of a simple closed
nearest-neighbor polygon on $\delta_n\Z^2$.  Its graph $G_n$ has every
lattice vertex in $\overline{D_n}$ and every nearest-neighbor edge whose
closed segment is contained in $\overline{D_n}$.  In particular, boundary
edges are part of $G_n$.

For two marked boundary points $a,b$, we require counterclockwise
parametrizations $\beta_n,\beta:[0,1]\to\C$, injective on $[0,1)$,
such that
\begin{gather}\label{intro:marked-jordan}
 \beta_n(0)=\beta_n(1)=a_n,\qquad \beta_n(1/2)=b_n,\qquad
 \beta(0)=\beta(1)=a,\qquad \beta(1/2)=b,\\
 \sup_{t\in[0,1]}|\beta_n(t)-\beta(t)|\longrightarrow0.\notag
\end{gather}
The free arc $F_n$ goes counterclockwise from $a_n$ to $b_n$; the other
arc $W_n$ goes counterclockwise from $b_n$ to $a_n$.  Wire together all
vertices of $W_n$, including the two marked endpoints.  All other
boundary vertices remain singleton blocks.  Remove from the sampled edge
set the boundary edges of $W_n$ and declare those edges open for the
interface drawing.  This is equivalent on all remaining edges to
\eqref{intro:finite-law} with the same wired partition: the endpoints of
each removed edge were already identified.  Edges on $F_n$ remain random.

The medial graph has vertices at the midpoints of primal edges, with
medial edges joining the midpoints of consecutive sides of a lattice
square.  At an open primal edge the four incident medial half-edges are
paired in quarter turns around the adjacent dual faces; at a closed
primal edge they are paired around the primal endpoints.  At the
boundary one uses the turns between the primal wired arc and its
complementary dual wired arc.  There is a unique medial strand between
the two changes of boundary condition.  We orient it from $a_n$ to $b_n$
and denote it by $\eta_n$.  This is the usual Dobrushin construction
\citep[Section~5.2.2]{DC2017Lectures}.  Occurrence order at a twice-visited
medial vertex is part of the curve.  The tile drawings used later preserve
this order and differ from the straight medial drawing by at most a
constant times $\delta_n$ in uniform curve distance.

For continuous oriented curves $\gamma_1,\gamma_2:[0,1]\to\C$, write
\begin{equation}\label{intro:curve-metric}
 d_{\mathrm{curv}}(\gamma_1,\gamma_2)
 =\inf_{\varphi_1,\varphi_2}
   \sup_{0\le t\le1}
   |\gamma_1(\varphi_1(t))-\gamma_2(\varphi_2(t))|,
\end{equation}
where the infimum is over increasing homeomorphisms of $[0,1]$.
Zero-distance curves are identified.  Chordal $\SLE_\kappa$ in $(D;a,b)$
is the conformal image of the upper-half-plane Loewner trace driven by
$\sqrt\kappa B_t$ from $0$ to $\infty$, with $B$ a standard Brownian
motion.  The remaining positive dilation of the conformal coordinate
does not change this law.

\subsection{Complete loop collections}
\label{intro:collections}

Let $\phi_q^0$ be the weak limit of \eqref{intro:finite-law} in increasing
square boxes of $\Z^2$ with free boundary conditions.  Form its complete
medial interface collection, using the same local pairings as above, and
then scale the embedding by $\delta$.  Denote the resulting collection
by $\Gamma_\delta$.  It contains every loop at every nesting level.
The elementary circuit argument in Section~\ref{sec:preliminaries}
shows that its discrete interfaces are almost surely finite.

Individual loops are unrooted and unoriented continuous maps of the
circle into the Riemann sphere $\widehat\C$.  Their distance is the
infimum of the uniform spherical distance over circular
orientation-preserving homeomorphisms and reversal, with zero-distance
representatives identified.  We retain distinct loops even when their
images coincide.  For two collections, an $\eps$-matching is a partial
bijection that covers all loops of spherical diameter greater than
$\eps$ in both collections, has matched-loop distance at most $\eps$,
and leaves only loops of diameter at most $\eps$ unmatched.  Convergence
of collections means the existence of such matchings with
$\eps\downarrow0$.  The continuum law in the next statement is the
standard nested whole-plane $\CLE_\kappa$ in this topology.

\begin{theorem}[Interfaces and all nested loops]
\label{intro:main}
For every fixed real $1\le q<4$, the following assertions hold.
\begin{enumerate}[label=\textup{(\roman*)}]
\item For every bounded Jordan domain $D$ with distinct boundary marks
$a,b$, and every approximation satisfying
\eqref{intro:marked-jordan}, the Dobrushin strand $\eta_n$ converges in
law, in \eqref{intro:curve-metric}, to chordal $\SLE_{\kappa(q)}$ from
$a$ to $b$ in $D$.
\item Under the specified free infinite-volume law,
$\Gamma_\delta$ converges in law as $\delta\downarrow0$ to the nested
whole-plane $\CLE_{\kappa(q)}$, in the spherical collection topology
just defined.  The limit is M\"obius invariant.
\end{enumerate}
Both conclusions hold along every mesh sequence.  No uniformity in $q$
is asserted.
\end{theorem}

The word ``critical'' refers to the edge parameter $p_q$, whose
criticality is established in \citet{BeffaraDuminilCopin2012}.
This gives a positive resolution, for $1\le q<4$, of the FK-to-SLE
prediction in \citet[Conjecture~9.7]{RohdeSchramm2005} and
\citet[Problem~2.6]{Schramm2007Problems}.  Those formulations also include
$0<q<1$; the result proved here is the stated parameter-restricted part.
The Jordan-domain theorem includes the square-domain formulation of
\citet{RohdeSchramm2005}.
The parameter $q=4$ is outside the scope of this article.  The strict
angular and six-arm estimates in the proof are used only at fixed
$q<4$.

\subsection{Square-lattice bond crossings}
\label{intro:crossings}

For four distinct counterclockwise boundary marks $a,b,c,d$ of $D$,
require the same uniform convergence as in
\eqref{intro:marked-jordan}, with the respective four marks at
parameters $0,1/4,1/2,3/4$.  Sample \emph{all} edges of $G_n$
independently with probability $1/2$, including every boundary edge, and
make no identifications or external connections.  Let $C_n$ be the
event of an actual open path from the closed counterclockwise arc
$[a_n,b_n]$ to the closed counterclockwise arc $[c_n,d_n]$.

\begin{theorem}[Cardy's formula for square-lattice bond percolation]
\label{intro:cardy}
Let $\psi:D\to\H$ send $a,c,d$ to $0,1,\infty$, and put
$x=\psi(b)\in(0,1)$.  Under the boundary and approximation conventions
above,
\begin{equation}\label{intro:cardy-formula}
 \lim_{n\to\infty}\mathbb P(C_n)
 =\frac{\displaystyle\int_0^x[t(1-t)]^{-2/3}\,\dd t}
        {\displaystyle\int_0^1[t(1-t)]^{-2/3}\,\dd t}.
\end{equation}
\end{theorem}

Theorem~\ref{intro:cardy} proves the square-lattice bond-percolation form
of Cardy's prediction with the free-boundary convention specified here.
We obtain it directly from the alternating four-change calculation and
a deterministic comparison on a common lattice graph.  In particular,
its proof supplies the boundary transfer for the actual crossing event;
it does not infer continuity of that event from a curve limit.

The same Cardy function also governs limiting annealed closed-cell crossing
probabilities in critical Poisson--Voronoi percolation in every fixed
bounded Jordan quadrilateral~\citep[Theorem~1.1]{OpenAIVoronoiCardy2026}.

\subsection{Square-lattice bond critical exponents}
\label{intro:percolation-exponents}

Write $\mathbb P_p^\square$ for independent bond percolation on $\Z^2$,
and put $B_R=\{z\in\C:|z|<R\}$.  Let
$\pi_1(R)=\mathbb P_{1/2}^\square(0\leftrightarrow\Z^2\setminus B_R)$.
For $k\ge1$, let $\pi_{2k}(r,R)$ be the $\mathbb P_{1/2}^\square$-probability
that $2k$ disjoint
actual primal-open and dual-open paths, cyclically alternating in
color, join $\overline B_r$ to $\C\setminus B_R$.  Here a dual-open
edge crosses a closed primal edge.  For the open cluster $C(0)$, set
\begin{align*}
 \theta(p)&=\mathbb P_p^\square(|C(0)|=\infty),\\
 \chi(p)&=\mathbb E_p^\square[|C(0)|;\ |C(0)|<\infty],\\
 \xi(p)^2&=\frac{1}{\chi(p)}
   \sum_{x\in\Z^2}|x|^2
   \mathbb P_p^\square(0\leftrightarrow x,\ |C(0)|<\infty)
   \qquad(p\ne\tfrac12).
\end{align*}
Thus $\chi$ is the truncated susceptibility and $\xi$ is the
second-moment correlation length.

\begin{corollary}[Square-bond arm and thermodynamic exponents]
\label{intro:bond-exponents}
For every fixed $k\ge1$ and every sufficiently large fixed $r$,
\begin{equation}\label{intro:bond-arm-exponents}
 \pi_1(R)=R^{-5/48+o(1)},\qquad
 \pi_{2k}(r,R)=R^{-(4k^2-1)/12+o(1)}
 \qquad(R\to\infty).
\end{equation}
Moreover,
\begin{equation}\label{intro:bond-thermodynamic-exponents}
\begin{aligned}
 \theta(p)&=(p-\tfrac12)^{5/36+o(1)}
   &&(p\downarrow\tfrac12),\\
 \chi(p)&=|p-\tfrac12|^{-43/18+o(1)}
   &&(p\to\tfrac12,\ p\ne\tfrac12),\\
 \xi(p)&=|p-\tfrac12|^{-4/3+o(1)}
   &&(p\to\tfrac12,\ p\ne\tfrac12).
\end{aligned}
\end{equation}
\end{corollary}

Section~\ref{sec:percolation-exponents} proves this consequence by
buffered annular comparisons using the complete loop limit, followed
by square-bond arm quasi-multiplicativity and the near-critical
scaling relations.  The arm assertion is for even, cyclically
alternating words; it makes no assertion for odd or other color words.
None of these exponent claims is made for $q\ne1$.

\subsection{Proof structure and technical contributions}

For $1\le q<4$, use the parameters
\begin{equation}\label{intro:lambda-rho}
 d=\sqrt q=2\cos\lambda,\qquad
 \rho=\lambda/\pi,\qquad 0<\lambda\le\pi/3.
\end{equation}
The first part of the proof determines an alternating-boundary
connection probability.  Four marked points divide the boundary into
arcs $A,B,C,D$, with the first and third assigned one color and
the other two the opposite color.  We study the event of an actual
path of the first color from $A$ to $C$.  Section~\ref{sec:fourchange} defines the
auxiliary finite law with four unmatched medial ends and no extra
pairing weights, and computes its limiting probability $f(\chi)$ when
the boundary marks have half-plane coordinates $0,\chi,1,\infty$,
$0<\chi<1$.  Completing the two designated arcs by separate
deterministic wires gives instead $f/(f+d(1-f))$.  We keep this
normalization change explicit before returning to actual FK crossings.

\paragraph{Positive comparisons before complex observables.}
Section~\ref{sec:preliminaries} identifies the capped tile law exactly
with the random-cluster law and establishes the collar and passage-count
comparisons used later.  It keeps actual components distinct from
boundary partition blocks.  Bulk box crossing at $q\le4$ and the stronger
up-to-boundary crossing estimates for $q<4$ have different roles
\citep{DCLM2019Isoradial,DMT2021Fractal}.  Section~\ref{sec:exponents}
derives the local exponent inequalities needed to control corner, tip,
and source factors.  All complex currents are subsequently estimated in
these positive laws.

\paragraph{Angular transfer and boundary holomorphicity.}
The four-change construction gives a discrete complex field depending
on an interior probe; its boundary values can be computed exactly.
To identify its limit with the conformal boundary map, we must prove
that every subsequential limit $h$ is holomorphic.  The finite row and
path calculations in Sections~\ref{sec:angular} and \ref{sec:current}
supply an angular decay estimate.  We apply it to two copies of the
field placed on opposite sides of a wall, one probed inside and the
other outside.  Repeated reflections keep the external states in
their norms, and the logarithmic estimates pay their complete costs.

The resulting mixed difference estimate tends to zero after division
by the mesh squared.  Closing the finitely many wall gaps and separating
the distant exterior copy yields
$\bar\partial h(z)\,\partial_jJ(u)=0$ in distributions, where $J$ is
the limiting exterior field and $\partial_j$ is a coordinate derivative.
The prescribed boundary values make $J$ nonconstant, so a suitable
exterior test forces $\bar\partial h=0$.
Sections~\ref{sec:pairedwall}--\ref{sec:holomorphic} prove these steps,
including control of the comparison factors through the successive limits.
The boundary-value calculation of Section~\ref{sec:fourchange} then
identifies the connection probability.

\paragraph{From fixed polygons to marked Jordan domains.}
The first calculation is made in rational orthogonal polygons with
specified lattice corner parities.  Proposition~\ref{jordan:transfer}
in Section~\ref{sec:jordan} extends it to arbitrary uniform marked
Jordan approximations.  Easier and harder comparison quads are drawn
on the same lattice and compared on their common edges.  Monotonicity
and a deterministic
comparison of actual paths sandwich its crossing probability between
their limits.  This step also proves the $q=1$ Cardy statement with all
boundary edges sampled freely.

\paragraph{Moving domains and the chordal driver.}
The fixed-shape formula must remain useful after a random exploration.
Section~\ref{sec:movingtests} proves this transfer using a finite chart
approximation and an averaged six-arm estimate under the original
switch law.  It does not require a uniform six-arm theorem in every
possible explored domain.  Section~\ref{sec:chordal} first normalizes a
shrinking test into a bounded martingale and then transfers it to the
true Dobrushin law.  Two tests identify the driver and its bracket.
Tightness and a separate traversal argument give the Euclidean uniform
curve conclusion of Theorem~\ref{intro:main}(i) for $q<4$.

\paragraph{Joint branches and loop identities.}
Sections~\ref{sec:peeling}--\ref{sec:loops} construct a target tree by
uniform-walled peeling.  They establish the reflected angular dynamics,
exclude forcepoint variation supported on endpoint times, and identify
independent continuation noises after separation, conditional on the
common causal history at separation.  The resulting branching SLE description is then matched
to the original discrete loops: winding information determines loop
labels, while additional estimates determine full traversal order.
An escape bound upgrades convergence in bounded windows to the
spherical matching topology for every nesting generation.

\subsection{Conventions on constants and limits}

Every theorem is for a fixed $q$.  Constants may depend on $q$,
macroscopic shapes, positive separations, and an explicitly fixed
finite set of tests.  Mesh limits are taken with those data fixed;
auxiliary separations, gaps, cutoffs, and localization parameters are
removed in the stated order.  A claim uniform over random exploration
histories is made only where the corresponding comparison has been
proved.  In particular, convergence in probability over histories is
not replaced by a deterministic estimate for every history.

The proof uses established public results at their stated hypotheses.
It uses no theorem about an annealed random planar map, no random
embedding, and no transfer from a random-map metric limit to the square
lattice.  The finite Dobrushin experiments and the free infinite-volume
experiment are different probability spaces; the theorems assert their
separate convergence and require no coupling between their samples.
